%% ch02-background.tex — Chapter II: Background and Related Work
%% Shared foundation for Chapters III–VIII. Section labels sec:bg:* are
%% referenced by other chapters; do not rename them.
\chapter{Background and Related Work}\label{ch:background}

This chapter provides the shared technical foundation for the research chapters that follow and situates this dissertation within the literature. The thesis statement of \cref{ch:intro} holds that useful quantum computation in the coming decade will be delivered by a co-designed stack in which classical high-performance computing (HPC) surrounds the quantum processing unit (QPU) at every layer. \Cref{fig:bg:stack} previews that stack and indicates which research question (RQ1--RQ5) each layer serves; the sections of this chapter walk through the layers in turn. We begin with the elements of gate-model quantum computation and the native gate sets of the superconducting and trapped-ion processors used throughout (\cref{sec:bg:qc}), then review the classical simulation of quantum circuits and the HPC systems on which it runs (\cref{sec:bg:sim}). We then treat the noise of near-term hardware and its mitigation (\cref{sec:bg:nisq}), variational optimization and the quantum approximate optimization algorithm (\cref{sec:bg:vqa}), classical-to-quantum data encoding (\cref{sec:bg:encoding}), quantum machine learning (\cref{sec:bg:qml}), circuit cutting and knitting (\cref{sec:bg:cutting}), quantum error correction (\cref{sec:bg:qec}), the cryptographic consequences of quantum computing and the post-quantum standards that respond to them (\cref{sec:bg:crypto}), and large language models for code synthesis (\cref{sec:bg:llm}). \Cref{sec:bg:summary} closes by identifying, for each research question, the gap in the literature that the dissertation addresses. Textbook material follows Nielsen and Chuang \cite{nielsen2010qcqi}; the research chapters cross-reference the sections of this chapter rather than repeating them.

% alt: Block diagram of the HPC-surrounded QPU stack. A central QPU box is surrounded by eight classical-HPC boxes (problem and data, GPU simulation, LLM/RL synthesis, encoding and cutting, structured optimizers, error mitigation, QEC decoding and audit, PQC module), each tagged with the research question it serves, all enclosed in a dashed frame labelled classical HPC.
\begin{figure}[htbp]
\centering
\begin{tikzpicture}[
  font=\small,
  box/.style={draw, rounded corners=2pt, align=center, inner sep=3pt, text width=36mm, minimum height=13mm, fill=cbSky!15},
  core/.style={box, fill=cbOrange!30, very thick, minimum height=18mm},
  arr/.style={-{Stealth[length=2.2mm]}, thick, draw=cbGray},
  bi/.style={{Stealth[length=2.2mm]}-{Stealth[length=2.2mm]}, thick, draw=cbGray}]
\node[box]  (data) at (-5.1, 3.0) {Problem and data\\ {\scriptsize QUBO, images, PDE fields}};
\node[box]  (sim)  at ( 0.0, 3.0) {GPU simulation\\ {\scriptsize \cudaq{} kernels (RQ1)}};
\node[box]  (llm)  at ( 5.1, 3.0) {LLM/RL synthesis\\ {\scriptsize rubric reward (RQ4)}};
\node[box]  (enc)  at (-5.1, 0.0) {Encoding and cutting\\ {\scriptsize \qcrank{}, knitting (RQ3)}};
\node[core] (qpu)  at ( 0.0, 0.0) {\textbf{QPU}\\ {\scriptsize IBM Heron, IonQ Forte}};
\node[box]  (opt)  at ( 5.1, 0.0) {Structured optimizers\\ {\scriptsize direct angles, walks (RQ2)}};
\node[box]  (mit)  at (-5.1,-3.0) {Error mitigation\\ {\scriptsize ZNE, readout (RQ2)}};
\node[box]  (qec)  at ( 0.0,-3.0) {QEC and audit\\ {\scriptsize seeded encoders (RQ5)}};
\node[box]  (pqc)  at ( 5.1,-3.0) {PQC module\\ {\scriptsize FIPS 140-3 (RQ5)}};
\draw[arr] (data) -- (enc);
\draw[arr] (data) -- (sim);
\draw[arr] (llm) -- (sim) node[midway, above, font=\scriptsize] {circuits};
\draw[bi]  (sim) -- (qpu) node[midway, right, font=\scriptsize] {verify, calibrate};
\draw[arr] (enc) -- (qpu);
\draw[bi]  (opt) -- (qpu);
\draw[arr] (qpu) -- (qec) node[midway, right, font=\scriptsize] {syndromes};
\draw[arr] (qpu.south west) -- (mit.north east);
\draw[arr, dashed] (qec) -- (pqc) node[midway, above, font=\scriptsize] {standards};
\begin{scope}[on background layer]
\node[draw=cbGray, dashed, rounded corners=4pt, fit=(data)(llm)(mit)(pqc), inner sep=7pt] (frame) {};
\end{scope}
\node[anchor=south west, font=\footnotesize\itshape, inner sep=1pt] at (frame.north west) {Classical HPC: Perlmutter A100 GPUs, Slurm, Podman-HPC/Shifter containers};
\end{tikzpicture}
\caption{The HPC-surrounded QPU stack studied in this dissertation. Every layer around the processor is classical and runs on HPC resources; the label in parentheses gives the research question that the layer serves. Schematic.}
\label{fig:bg:stack}
\end{figure}

\section{Quantum Computation Fundamentals}\label{sec:bg:qc}

\subsection{Qubits, Gates and Measurement}

A qubit is a two-level quantum system whose pure states are unit vectors in $\C^2$, written $\ket{\psi}=\alpha\ket{0}+\beta\ket{1}$ with $|\alpha|^2+|\beta|^2=1$. An $n$-qubit register lives in the tensor-product space $(\C^2)^{\otimes n}\cong\C^{2^n}$, and a general pure state is
\begin{equation}\label{eq:bg:statevector}
\ket{\psi}=\sum_{x\in\{0,1\}^n}\psi_x\ket{x},\qquad \sum_{x}|\psi_x|^2=1,
\end{equation}
where the $2^n$ complex amplitudes $\psi_x$ are indexed by bit strings. The exponential size of this description is at once the source of the hoped-for computational power and, as \cref{sec:bg:sim} discusses, the reason classical simulation runs into a memory wall. Mixed states are described by density operators $\rho$, positive semidefinite with unit trace, and noise processes act as completely positive trace-preserving maps $\rho\mapsto\sum_k E_k\rho E_k^\dagger$ with Kraus operators satisfying $\sum_k E_k^\dagger E_k=I$ \cite{nielsen2010qcqi}.

Closed-system evolution is unitary, and a quantum circuit is a sequence of unitary gates acting on one or a few qubits at a time. The single-qubit Pauli matrices $X$, $Y$ and $Z$, the Hadamard gate $H$, the phase gates $S=\mathrm{diag}(1,i)$ and $T=\mathrm{diag}(1,e^{i\pi/4})$, and the rotations $R_P(\theta)=e^{-i\theta P/2}$ for $P\in\{X,Y,Z\}$ appear throughout this dissertation, together with the two-qubit controlled-NOT (CNOT or CX) and controlled-$Z$ (CZ) gates and the parameterized interaction $e^{-i\theta Z_iZ_j/2}$, which is the building block of the QAOA circuits of \cref{sec:bg:vqa}. A computation ends with measurement in the computational basis, which returns the bit string $x$ with probability $|\psi_x|^2$ and collapses the state. Because one execution yields one sample, quantities of interest such as the expectation value $\expval{O}{\psi}$ of an observable $O$ must be estimated from $S$ repeated executions, or shots, with a statistical error that scales as $1/\sqrt{S}$. This shot cost is invisible in idealized complexity analyses and is a recurring theme of \cref{ch:qml,ch:synthesis}.

\subsection{Universal and Native Gate Sets}

A gate set is universal if every unitary on $n$ qubits can be approximated to arbitrary precision by circuits over it. The Clifford$+T$ set $\{H,S,T,\mathrm{CNOT}\}$ is universal and dominates fault-tolerant compilation (\cref{sec:bg:qec}), and the Solovay--Kitaev theorem guarantees that approximating a single-qubit gate to accuracy $\eps$ requires only $\mathrm{polylog}(1/\eps)$ gates from a finite universal set \cite{nielsen2010qcqi}. Physical devices, however, implement a small \emph{native} gate set dictated by their control hardware, and every algorithmic circuit must be rewritten in that set before it can run. IBM's fixed-frequency transmon processors of the Falcon and Eagle generations used the cross-resonance interaction to realize an entangling gate; the echoed cross-resonance (ECR) sequence cancels the dominant unwanted terms of that interaction and served as the native two-qubit gate of the 127-qubit Eagle devices \cite{sheldon2016echoedcr,bravyi2022futuresc}. The Heron generation used in this dissertation replaces the cross-resonance gate with a CZ gate mediated by tunable couplers, complemented by the single-qubit $\sqrt{X}$, $X$ and virtual $R_Z$ gates, the last of which are implemented as frame changes at zero duration and zero error \cite{mckay2017virtualz}\todo{verify the Heron native gate set (CZ, $\sqrt{X}$, X, RZ) against the backend properties of each device in \cref{tab:bg:devices}}. Trapped-ion processors such as IonQ's Aria and Forte systems use the M{\o}lmer--S{\o}rensen interaction, which couples pairs of ions through a shared motional mode and produces an $XX$-type entangling gate between arbitrary pairs, giving all-to-all connectivity at the expense of slower gates \cite{sorensen1999ms}. Native gate sets matter because the two-qubit gate count after rewriting, not the abstract gate count of the algorithm, determines the error budget of a NISQ execution.

\subsection{Transpilation and Qubit Routing}

Transpilation converts an abstract circuit into one that respects a device's native gate set and coupling graph. It comprises gate decomposition, layout (the assignment of logical to physical qubits), routing (the insertion of SWAP gates so that every two-qubit gate acts on physically coupled qubits) and gate-level optimization such as the cancellation of adjacent inverse gates and the fusion of single-qubit rotations. Routing is NP-hard in general; the SABRE heuristic of Li et al. \cite{li2019sabre}, which \qiskit{} uses by default, performs a bidirectional search that greedily minimizes a distance-based cost while inserting SWAPs \cite{javadiabhari2024qiskit}. Since each SWAP costs three CNOTs (or CZs with additional single-qubit gates) on a heavy-hex lattice, routing can double or triple the two-qubit depth of a circuit whose logical connectivity does not match the hardware. This is why the algorithm designs of \cref{ch:nisq,ch:qml} pay close attention to hardware-native structure, and why the hardware-aware cut selection of \cref{ch:qml} takes calibration data and physical qubit distance as first-class inputs.

\section{Classical Simulation of Quantum Circuits at HPC Scale}\label{sec:bg:sim}

\subsection{State-Vector Simulation and the Memory Wall}

The most direct classical simulator stores all $2^n$ amplitudes of \cref{eq:bg:statevector} and applies each gate as a sparse linear map. A single-qubit gate $U=(u_{ab})$ on target qubit $t$ acts on every pair of amplitudes whose indices differ only in bit $t$,
\begin{equation}\label{eq:bg:gateapply}
\begin{pmatrix}\psi'_{x\,[t\leftarrow 0]}\\ \psi'_{x\,[t\leftarrow 1]}\end{pmatrix}
=\begin{pmatrix}u_{00}&u_{01}\\ u_{10}&u_{11}\end{pmatrix}
\begin{pmatrix}\psi_{x\,[t\leftarrow 0]}\\ \psi_{x\,[t\leftarrow 1]}\end{pmatrix}
\qquad\text{for all $x$ with bit $t$ cleared},
\end{equation}
so that one gate costs $\Theta(2^n)$ arithmetic operations and one sweep over memory, and a $k$-qubit gate, which touches $2^k$-tuples of amplitudes, costs $\Theta(2^{n+k})$. A circuit of $G$ gates therefore costs $\Theta(G\,2^n)$ time, but time is rarely the binding constraint. The state vector itself occupies
\begin{equation}\label{eq:bg:memory}
M(n)=2^n\times 16\ \text{bytes}\qquad(\text{double-precision complex amplitudes}),
\end{equation}
or half that in single precision. \Cref{fig:bg:memwall} plots \cref{eq:bg:memory} against the memory of the NVIDIA A100 accelerators \cite{choquette2021a100} used throughout this dissertation. The curve doubles with every qubit: a 30-qubit state occupies \SI{17}{\giga\byte}, a 40-qubit state \SI{17.6}{\tera\byte}. The 32-qubit single-GPU, 34-qubit four-GPU and 42-qubit, 1024-GPU simulations reported for \qgear{} \cite{guo2025qgear} sit exactly where single-precision amplitudes fill the available device memory, which shows that at scale state-vector simulation is a memory-capacity problem first and an arithmetic problem second. This is the memory wall behind RQ1.

% alt: Log-scale line plot of state-vector memory in gigabytes versus qubit count from 24 to 46, one solid line for double precision and one dashed line for single precision, with dotted and dash-dotted horizontal reference lines for a 40 GB A100, an 80 GB A100, a four-GPU node and 1024 GPUs, and three diamond markers at 32, 34 and 42 qubits on the single-precision line.
\begin{figure}[htbp]
\centering
\begin{tikzpicture}
\begin{axis}[
  width=0.88\linewidth, height=0.72\linewidth,
  ymode=log, xmin=24, xmax=46, ymin=0.05, ymax=2e6,
  xlabel={Number of qubits $n$}, ylabel={State-vector memory (GB)},
  xtick={24,28,...,44},
  legend pos=south east]
\addplot[cbBlue, thick, solid, mark=none] coordinates {(24,0.2684) (26,1.074) (28,4.295) (30,17.18) (32,68.72) (34,274.9) (36,1100) (38,4398) (40,1.759e+04) (42,7.037e+04) (44,2.815e+05) (46,1.126e+06)};
\addlegendentry{double precision, 16 B per amplitude}
\addplot[cbOrange, thick, dashed, mark=none] coordinates {(24,0.1342) (26,0.5369) (28,2.147) (30,8.59) (32,34.36) (34,137.4) (36,549.8) (38,2199) (40,8796) (42,3.518e+04) (44,1.407e+05) (46,5.629e+05)};
\addlegendentry{single precision, 8 B per amplitude}
\addplot[cbRed, only marks, mark=diamond*, mark size=3.2pt] coordinates {(32,34.36) (34,137.4) (42,3.518e4)};
\addlegendentry{\qgear{} runs on 1, 4 and 1024 GPUs}
\addplot[cbGray, dotted, thick, mark=none, forget plot] coordinates {(24,40) (46,40)} node[pos=1, below left, font=\scriptsize, inner sep=1pt] {one A100, \SI{40}{\giga\byte}};
\addplot[cbGray, dotted, thick, mark=none, forget plot] coordinates {(24,80) (46,80)} node[pos=0.02, above right, font=\scriptsize, inner sep=1pt] {one A100, \SI{80}{\giga\byte}};
\addplot[cbGray, dashdotted, thick, mark=none, forget plot] coordinates {(24,160) (46,160)} node[pos=0.02, above right, font=\scriptsize, inner sep=1pt] {one node, four \SI{40}{\giga\byte} GPUs};
\addplot[cbGray, dashdotted, thick, mark=none, forget plot] coordinates {(24,40960) (46,40960)} node[pos=0.02, above right, font=\scriptsize, inner sep=1pt] {1024 GPUs, \SI{40.96}{\tera\byte} in total};
\end{axis}
\end{tikzpicture}
\caption{State-vector memory footprint $2^n\times 16$ bytes (double precision) and $2^n\times 8$ bytes (single precision) as a function of the qubit count, compared with the memory of one NVIDIA A100 (40 or \SI{80}{\giga\byte}), one Perlmutter GPU node with four \SI{40}{\giga\byte} A100s, and 1024 such GPUs. Diamonds mark the 32-qubit single-GPU, 34-qubit four-GPU and 42-qubit 1024-GPU simulations reported in \cite{guo2025qgear}, placed on the single-precision curve. Schematic; reference lines are nominal device capacities.}
\label{fig:bg:memwall}
\end{figure}

\subsection{Distributed State Vectors on GPUs}

Passing the wall requires distributing the amplitudes across the memories of many devices. The standard scheme partitions on the high-order index bits: with $2^m$ devices, the $m$ most significant bits of $x$ select the device and the remaining $n-m$ bits address the local slice. Gates acting on any of the $n-m$ low-order, or local, qubits then need no communication, because both amplitudes of each pair in \cref{eq:bg:gateapply} reside on the same device. Gates on the $m$ high-order, or global, qubits pair amplitudes held by different devices and are executed either by pairwise exchange of half-slices or, more efficiently, by periodically permuting global and local qubits through an all-to-all exchange so that a block of subsequent gates becomes local \cite{bayraktar2023cuquantum}. On Perlmutter the exchange traverses NVLink within a node and the HPE Slingshot network between nodes through MPI, so that interconnect bandwidth rather than floating-point throughput bounds the achievable speed once the state exceeds one node. \qgear{} \cite{guo2025qgear} relies on the multi-GPU distribution of \cudaq{} and reported close to full GPU utilization on up to 1024 A100 GPUs.

\subsection{Tensor-Network Alternatives}

Not every state requires $2^n$ numbers. A matrix product state (MPS) writes the amplitudes as a product of matrices,
\begin{equation}\label{eq:bg:mps}
\psi_{x_1x_2\cdots x_n}=\sum_{\alpha_1,\dots,\alpha_{n-1}}A^{[1]x_1}_{\alpha_1}\,A^{[2]x_2}_{\alpha_1\alpha_2}\cdots A^{[n]x_n}_{\alpha_{n-1}},
\end{equation}
where each bond index $\alpha_i$ ranges over the bond dimension $\chi$, so that storage is $\bigO{n\chi^2}$ rather than $\bigO{2^n}$ \cite{perezgarcia2006mps}. Any state can be represented exactly if $\chi$ may grow to $2^{\lfloor n/2\rfloor}$, but the representation is efficient only when the entanglement across every cut of the chain is bounded, as in shallow circuits or circuits with limited long-range entanglement. Each entangling gate multiplies $\chi$ until it is truncated by a singular-value decomposition, which introduces a controlled approximation error. General tensor-network contraction underlies the cuTensorNet library of cuQuantum \cite{bayraktar2023cuquantum} and was the basis of the classical counter-arguments in the quantum-supremacy debate \cite{arute2019supremacy}. In this dissertation MPS compilation with a bounded bond dimension is used in \cref{ch:qml} to compile sub-circuits of data-encoding circuits beyond the reach of exact simulation, while the stabilizer simulators of \cref{sec:bg:qec} cover the complementary Clifford regime.

\subsection{GPU Simulators and the HPC Software Context}

\Cref{tab:bg:simulators} summarizes the simulators that recur in this dissertation. \qiskit{} Aer \cite{javadiabhari2024qiskit} provides state-vector, density-matrix, MPS and stabilizer methods with OpenMP threading, cuStateVec-backed GPU execution and batched multi-shot GPU techniques \cite{doi2023multishotgpu}. \cudaq{} \cite{kim2023cudaquantum} is NVIDIA's kernel-based programming model whose simulation targets are built on the cuQuantum SDK \cite{bayraktar2023cuquantum}: cuStateVec for state vectors, including multi-GPU and multi-node distribution, and cuTensorNet for tensor networks. PennyLane Lightning \cite{asadi2024pennylanelightning} offers CPU, Kokkos and cuStateVec back ends and is organized around adjoint differentiation for variational workloads. Qulacs \cite{suzuki2021qulacs} and QCLAB++ \cite{vanbeeumen2023qclabpp} are compact C++ simulators with single-GPU support, and Stim \cite{gidney2021stim} is a stabilizer simulator restricted to Clifford circuits. A practical difficulty is that these tools expose different programming models: a \qiskit{} circuit cannot be handed to a \cudaq{} kernel without being rewritten. Closing that gap without asking algorithm developers to abandon \qiskit{} is the software contribution of \cref{ch:qgear}.

\begin{table}[htbp]
\caption{Circuit simulators referenced in this dissertation, characterized by simulation method, parallelism and GPU support. Compiled from the cited software papers; capabilities refer to the versions described there.}
\label{tab:bg:simulators}
\centering\footnotesize
\begin{tabularx}{\linewidth}{@{}>{\raggedright\arraybackslash}p{2.6cm}>{\raggedright\arraybackslash}X>{\raggedright\arraybackslash}X>{\raggedright\arraybackslash}X>{\raggedright\arraybackslash}p{2.3cm}@{}}
\toprule
Simulator & Methods & Parallelism & GPU support & Programming model \\
\midrule
\qiskit{} Aer \cite{javadiabhari2024qiskit,doi2023multishotgpu} & state vector, density matrix, MPS, stabilizer & OpenMP threads; MPI with cuStateVec & NVIDIA via Thrust and cuStateVec; batched shots & \qiskit{} circuits (Python) \\
\addlinespace[3pt]
\cudaq{} with cuQuantum \cite{kim2023cudaquantum,bayraktar2023cuquantum} & state vector (cuStateVec), tensor network (cuTensorNet) & multi-GPU and multi-node through MPI & native NVIDIA & quantum kernels in C++ or Python \\
\addlinespace[3pt]
PennyLane Lightning \cite{asadi2024pennylanelightning} & state vector with adjoint gradients & OpenMP, Kokkos, MPI & lightning.gpu (cuStateVec), lightning.kokkos & PennyLane QNodes (Python) \\
\addlinespace[3pt]
Qulacs \cite{suzuki2021qulacs} & state vector, density matrix & OpenMP and SIMD & single GPU (CUDA) & C++ and Python \\
\addlinespace[3pt]
QCLAB++ \cite{vanbeeumen2023qclabpp} & state vector & OpenMP & single GPU (CUDA) & header-only C++ \\
\addlinespace[3pt]
Stim \cite{gidney2021stim} & stabilizer tableau, Pauli-frame sampling & SIMD bit packing & none (CPU only) & C++ and Python; Clifford circuits only \\
\bottomrule
\end{tabularx}
\end{table}

The simulations in this dissertation ran predominantly on Perlmutter at NERSC, whose GPU nodes each combine an AMD EPYC processor with four NVIDIA A100 GPUs of \SI{40}{\giga\byte} high-bandwidth memory (\SI{80}{\giga\byte} on a subset of nodes) connected by NVLink, with nodes linked by the HPE Slingshot interconnect; CPU baselines used dual-socket AMD EPYC 7763 nodes with 128 cores and \SI{512}{\giga\byte} of memory \cite{guo2025qgear,choquette2021a100}. Jobs are scheduled by Slurm \cite{jette2023slurm}, which allocates nodes, launches MPI ranks and enforces time and memory limits, so that a simulator must be packaged as a batch job. Software environments are delivered as containers: Shifter \cite{gerhardt2017shifter} converts Docker images into a form that can be launched at scale on compute nodes, and its successor Podman-HPC \cite{stephey2022podman} provides rootless, OCI-compatible containers with GPU and MPI hooks, so that a \cudaq{} or PyTorch installation pinned to specific library versions can be rebuilt identically later; the artifacts of \cref{app:repro} follow this practice. The broader vision of a quantum supercomputer, in which QPUs form one accelerator class among many inside such a facility, is articulated by Mohseni et al. \cite{mohseni2024quantumsupercomputer} and frames the stack of \cref{fig:bg:stack}.

\section{Noise on NISQ Hardware and Error Mitigation}\label{sec:bg:nisq}

\subsection{Noise Channels and Device Metrics}

Preskill's term noisy intermediate-scale quantum (NISQ) \cite{preskill2018nisq} designates processors with tens to a few hundred qubits whose errors are not corrected but must be tolerated. The dominant incoherent processes are energy relaxation, characterized by the time $T_1$ over which the excited-state population decays as $e^{-t/T_1}$, and dephasing, characterized by $T_2\le 2T_1$, over which the off-diagonal elements of the density matrix decay as $e^{-t/T_2}$. A qubit that idles for a time $t$ while other qubits are operated on therefore accumulates an error of order $t/T_1$ and $t/T_2$ even though no gate acts on it; this idling decoherence is why long, sparse circuits are not free and why delay-aware scheduling matters. Gate errors are quantified by randomized benchmarking \cite{magesan2012rb}; on current superconducting devices single-qubit errors are of order $10^{-4}$ and two-qubit errors of order $10^{-3}$, for instance a median CZ error of \num{1.46e-3} on ibm\_pittsburgh during the experiments of \cref{ch:qml} \cite{guo2025shardq}. Readout error, the probability that a qubit prepared in $\ket{0}$ or $\ket{1}$ is reported as the other value, is typically of order $10^{-2}$, an order of magnitude worse than gate error. Finally, crosstalk denotes unwanted coupling between simultaneously driven or neighbouring qubits, including residual $ZZ$ interactions with spectator qubits; it is not captured by isolated gate benchmarks and degrades circuits whose gates run in parallel \cite{sheldon2016echoedcr}. \Cref{tab:bg:devices} lists the processors used in this dissertation together with the values reported in the source publications; cells that those sources do not report are left blank\todo{verify: fill the blank qubit counts, native gates and median two-qubit errors of \cref{tab:bg:devices} from the calibration snapshots archived with each experiment}.

\begin{table}[htbp]
\caption{Quantum processors used in this dissertation. Qubit counts and the median CZ error are those reported in \cite{guo2025deal,guo2025shardq,guo2025qawa,guo2025vqt,guo2025qpie,guo2026nnqa,guo2026rubriq}; a dash marks a value not reported in those sources. Native two-qubit gates are CZ for the IBM Heron generation and the M{\o}lmer--S{\o}rensen (MS) interaction for IonQ.}
\label{tab:bg:devices}
\centering\footnotesize
\begin{tabular}{@{}llccll@{}}
\toprule
Device & Family & Qubits & Native 2q gate & Median 2q error & Used in \\
\midrule
ibm\_torino & IBM Heron r1 & 133 & CZ & --- & \cref{ch:nisq} \\
ibm\_marrakesh & IBM Heron r2 & 156 & CZ & --- & \cref{ch:nisq,ch:qml} \\
ibm\_kingston & IBM Heron & --- & CZ & --- & \cref{ch:qml} \\
ibm\_pittsburgh & IBM Heron r3 & 156 & CZ & \num{1.46e-3} & \cref{ch:nisq,ch:qml} \\
ibm\_boston & IBM Heron r3 & 156 & CZ & --- & \cref{ch:qml} \\
ibm\_miami & IBM Nighthawk & --- & --- & --- & \cref{ch:qml} \\
IonQ Aria-1 & trapped ion & --- & MS & --- & \cref{ch:qml} \\
IonQ Forte-1 & trapped ion & 36 & MS & --- & \cref{ch:qml,ch:synthesis} \\
\bottomrule
\end{tabular}
\end{table}

The compound effect of these processes is captured by a simple heuristic. If a circuit has $G_2$ two-qubit gates of error $\eps_2$, $G_1$ single-qubit gates of error $\eps_1$ and $n$ measured qubits with readout error $\eps_r$, its success probability is roughly $(1-\eps_2)^{G_2}(1-\eps_1)^{G_1}(1-\eps_r)^{n}$ before idling and crosstalk are counted. With $\eps_2\approx10^{-3}$ the two-qubit budget alone falls below $e^{-1}$ at about a thousand gates, and in practice coherent errors and crosstalk make the usable depth considerably smaller: the \deal{} experiments of \cref{ch:nisq} observed performance degrading beyond roughly 41--43 CZ gates on Heron processors \cite{guo2025deal}. Depth is the enemy on NISQ hardware, and every design in this dissertation is, in one way or another, a strategy for buying accuracy with less depth.

\subsection{Error Mitigation}

Error mitigation reduces the bias of expectation values estimated on noisy hardware at the cost of additional circuit executions, without the qubit overhead of error correction. Zero-noise extrapolation (ZNE), introduced independently by Temme et al. \cite{temme2017errormitigation} and Li and Benjamin \cite{li2017activeerror}, executes the circuit at several amplified noise levels $\lambda_1<\lambda_2<\dots<\lambda_m$, with $\lambda=1$ the native level, and extrapolates the measured expectation values $E(\lambda_i)$ to $\lambda=0$. Under the model $E(\lambda)=E^{\ast}+\sum_{k=1}^{m-1}a_k\lambda^k$, Richardson extrapolation gives
\begin{equation}\label{eq:bg:zne}
E^{\ast}\approx\sum_{i=1}^{m}\gamma_iE(\lambda_i),\qquad \sum_{i}\gamma_i=1,\qquad \sum_{i}\gamma_i\lambda_i^{k}=0\quad(k=1,\dots,m-1),
\end{equation}
which for two points reduces to $E^{\ast}\approx(\lambda_2E(\lambda_1)-\lambda_1E(\lambda_2))/(\lambda_2-\lambda_1)$. Kandala et al. \cite{kandala2019errormitigation} amplified noise by stretching gate pulses, which requires pulse-level access. Digital ZNE \cite{giurgicatiron2020digitalzne} amplifies noise at the gate level instead: unitary folding replaces a circuit $U$, or a subset of its gates, by $U(U^{\dagger}U)^{m}$, which is logically the identity but multiplies the gate noise by an odd factor $\lambda=2m+1$, and identity or delay insertion lengthens idle periods so that the coherence-limited part of the noise is scaled separately from the gate-limited part. The adaptive digital ZNE of \cref{ch:nisq} \cite{guo2025deal} belongs to the latter family: it inserts delay gates whose durations are chosen from device calibration data so that crosstalk and coherence errors are controlled rather than merely amplified. The variance of the extrapolated estimator grows with $\sum_i\gamma_i^2$, so ZNE trades bias for shots, and it fails when the noise is not a smooth function of $\lambda$, which is one reason the extrapolation must be validated against simulation.

Measurement errors are mitigated separately. The readout channel is modelled as a stochastic matrix $A$ with $A_{yx}=\Pr[\text{read }y\mid\text{prepared }x]$, calibrated from preparation circuits, after which $\hat{p}_{\mathrm{ideal}}=A^{-1}\hat{p}_{\mathrm{noisy}}$ or a constrained least-squares variant recovers the ideal distribution; because $A$ has $2^n\times2^n$ entries, scalable methods assume a tensor-product or few-qubit-correlated structure \cite{bravyi2021measurement}. Probabilistic error cancellation \cite{temme2017errormitigation} inverts the gate noise channel as a quasi-probability distribution and is unbiased, but its sampling overhead grows exponentially with the number of noisy gates, which makes ZNE and readout mitigation the practical choice at the depths considered here.

\section{Variational Algorithms and QAOA}\label{sec:bg:vqa}

\subsection{Variational Quantum Algorithms}

Variational quantum algorithms (VQAs) \cite{cerezo2021vqa} pair a parameterized quantum circuit $U(\bm{\theta})$ with a classical optimizer: the QPU prepares $\ket{\psi(\bm{\theta})}=U(\bm{\theta})\ket{0}^{\otimes n}$ and estimates a cost $C(\bm{\theta})=\expval{H}{\psi(\bm{\theta})}$ from shots, the classical side updates $\bm{\theta}$, and the loop repeats. The circuits are shallow by design, which suits NISQ hardware. The variational quantum eigensolver targets ground-state energies of molecular or spin Hamiltonians, the quantum approximate optimization algorithm (QAOA) targets combinatorial problems, and the quantum neural networks of \cref{sec:bg:qml} target supervised learning. All share the same failure modes: the number of cost evaluations demanded by the optimizer, each of which costs shots and queue time, and the flatness of the optimization landscape discussed below. Because many evaluations are performed with slightly different parameters, GPU simulation is the natural development environment, and this dissertation uses it both to validate variational circuits before hardware submission and, in \cref{ch:synthesis}, as the reward oracle of a learning loop.

\subsection{QUBO Problems and the QAOA Ansatz}

Quadratic unconstrained binary optimization (QUBO) asks for $\bm{x}^{\ast}=\argmin_{\bm{x}\in\{0,1\}^n}\bm{x}^{\mathsf{T}}Q\bm{x}$ for a symmetric matrix $Q\in\R^{n\times n}$. Glover et al. \cite{glover2018qubo} survey how constraints are folded into $Q$ through penalty terms, and Monta{\~n}ez-Barrera et al. \cite{montanezbarrera2024unbalanced} show how inequality constraints can be encoded without slack variables. MaxCut on a graph $(V,E)$ is the canonical example, with $Q$ built from the adjacency matrix; the travelling-salesman problem uses $n^2$ binary variables $x_{v,t}$ (city $v$ at position $t$) with one-hot penalties on rows and columns; the knapsack problem penalizes capacity violations; and portfolio optimization minimizes the risk $\bm{x}^{\mathsf{T}}\Sigma\bm{x}$ minus the expected return under a budget penalty. Substituting $x_i=(1-Z_i)/2$ converts the QUBO objective into a diagonal cost Hamiltonian
\begin{equation}\label{eq:bg:hcost}
\Hcost=\sum_{i<j}J_{ij}Z_iZ_j+\sum_{i}h_iZ_i+c,\qquad J_{ij}=\tfrac12Q_{ij},\quad h_i=-\tfrac12\sum_{j}Q_{ij},
\end{equation}
whose ground state encodes $\bm{x}^{\ast}$. QAOA \cite{farhi2014qaoa} prepares the trial state
\begin{equation}\label{eq:bg:qaoa}
\ket{\bm{\gamma},\bm{\beta}}=\prod_{k=1}^{p}e^{-i\beta_k\Hmix}\,e^{-i\gamma_k\Hcost}\ket{+}^{\otimes n},\qquad \Hmix=\sum_{i=1}^{n}X_i,
\end{equation}
with $2p$ angles, and minimizes $\expval{\Hcost}{\bm{\gamma},\bm{\beta}}$. Each cost layer is a product of commuting $ZZ$ and $Z$ rotations, with every $e^{-i\gamma Z_iZ_j}$ compiling to two CNOTs and one $R_Z(2\gamma)$, and each mixer layer is a product of $R_X(2\beta)$ gates; \cref{fig:bg:qaoa} shows one layer. As $p\to\infty$ the ansatz can reproduce adiabatic evolution and reach the exact optimum, but on hardware the two-qubit gate count $2p|E|$ plus routing overhead binds long before that. Weidenfeller et al. \cite{weidenfeller2022qaoascaling} found that SWAP overhead dominates QAOA on heavy-hex hardware for dense problem graphs; Sachdeva et al. \cite{sachdeva2024qaoa127} reported that a 127-qubit gate-model device could outperform annealers on structured instances; and Abbas et al. \cite{abbas2024challengesqopt} survey the field's open problems critically. Continuous-time quantum walks \cite{childs2009quantumwalk} offer an alternative view of the mixer as a walk on the hypercube and inspire the approximate-walk formulation of \cref{ch:nisq}.

% alt: Three-qubit quantum circuit showing one QAOA layer: Hadamard gates on all qubits, a ZZ rotation on qubits 0 and 1 decomposed into CNOT, RZ of two gamma and CNOT, a boxed ZZ rotation on qubits 1 and 2, then RX of two beta mixers on every qubit followed by measurements.
\begin{figure}[htbp]
\centering
\begin{quantikz}[row sep={0.75cm,between origins}, column sep=0.3cm]
\lstick{$q_0=\ket{0}$} & \gate{H} & \ctrl{1}\gategroup[3,steps=4,style={dashed,rounded corners,inner xsep=2pt},background,label style={label position=below,anchor=north,yshift=-0.2cm}]{cost layer $e^{-i\gamma_k\Hcost}$} & & \ctrl{1} & & \gate{R_X(2\beta_k)}\gategroup[3,steps=1,style={dashed,rounded corners,inner xsep=2pt},background,label style={label position=below,anchor=north,yshift=-0.2cm}]{mixer} & \meter{} \\
\lstick{$q_1=\ket{0}$} & \gate{H} & \targ{} & \gate{R_Z(2\gamma_k)} & \targ{} & \gate[2]{e^{-i\gamma_k Z_1Z_2}} & \gate{R_X(2\beta_k)} & \meter{} \\
\lstick{$q_2=\ket{0}$} & \gate{H} & & & & & \gate{R_X(2\beta_k)} & \meter{}
\end{quantikz}
\caption{One QAOA layer for a three-qubit cost Hamiltonian with couplings on qubit pairs $(0,1)$ and $(1,2)$. The first $ZZ$ rotation is shown in its native decomposition into two CNOTs and one $R_Z(2\gamma_k)$; the second is drawn as a block. Single-qubit $Z_i$ terms of \cref{eq:bg:hcost} would add $R_Z$ gates to the cost layer. Schematic.}
\label{fig:bg:qaoa}
\end{figure}

\subsection{Trainability and Parameter Initialization}

The classical optimization of $\bm{\gamma}$ and $\bm{\beta}$ is the Achilles' heel of QAOA. McClean et al. \cite{mcclean2018barren} showed that for sufficiently random parameterized circuits, which approximate unitary 2-designs, the variance of every gradient component vanishes exponentially in $n$, $\mathrm{Var}[\partial_{\theta}C]\in\bigO{b^{-n}}$ for some $b>1$, so that gradients are exponentially small almost everywhere: the barren plateau. Holmes et al. \cite{holmes2022expressibility} connected the severity of the plateau to ansatz expressibility, and Wang et al. \cite{wang2021noisebp} showed that local noise induces a plateau whose severity grows with circuit depth independently of the ansatz, tying trainability directly to the depth budget of \cref{sec:bg:nisq}. Each gradient component moreover costs shots, and optimizer iterations multiply the hardware queue time. The parameter-initialization literature responds by exploiting structure. Zhou et al. \cite{zhou2020qaoa} observed that optimal QAOA angles concentrate and vary smoothly with $p$, permitting interpolation-based initialization from smaller $p$; Egger et al. \cite{egger2021warmstart} warm-start QAOA from the continuous relaxation of the QUBO, biasing the initial state and mixer toward a classical solution; and fixed-angle and parameter-transfer approaches reuse angles across instances of the same problem class. All of these still treat the angles as free parameters to be tuned. The \deal{} approach of \cref{ch:nisq} \cite{guo2025deal} goes one step further and maps the QUBO coefficients directly onto ansatz angles, so that optimization is replaced by a deterministic construction, and the \qawa{} method \cite{guo2025qawa} replaces the black-box loop with a shallow, classically traceable oracle whose mid-circuit measurement statistics are interpretable without full state tomography.

\section{Encoding Classical Data into Quantum States}\label{sec:bg:encoding}

\subsection{Basis, Angle and Amplitude Encodings}

Any quantum algorithm that processes classical data must first load it into a quantum state, and the choice of encoding determines both the circuit cost and what can be computed afterwards \cite{schuld2021dataencoding,larose2020encodings}. Basis encoding writes a bit string into the computational basis, $\bm{b}\mapsto\ket{\bm{b}}$; it is exact and cheap but uses one qubit per bit. Angle encoding maps each real feature $x_j$ to a rotation of one qubit, $\ket{\psi(\bm{x})}=\bigotimes_j R_Y(x_j)\ket{0}$, so that $n$ features cost $n$ qubits and constant depth; repeated or higher-frequency angle encodings act as feature maps whose Fourier spectrum controls the expressivity of the resulting model \cite{schuld2021dataencoding,havlicek2019featurespaces}. Amplitude encoding stores $N=2^n$ normalized values as the amplitudes $\psi_x$ of an $n$-qubit state and achieves exponential compression in qubits, but preparing an arbitrary amplitude-encoded state requires $\Theta(2^n)$ gates in general and reading the values back requires tomography, so the compression is paid for at both ends. Proposals that assume amplitude-encoded inputs at no cost, including the quantum linear-systems algorithm \cite{harrow2009hhl}, are exposed to this input problem \cite{biamonte2017qml}.

\subsection{QPIXL and QCrank}

Image and array data motivated a family of encodings that keep the exponential compression of amplitude encoding for the \emph{address} but store the \emph{values} as rotation angles. The QPIXL framework of Amankwah et al. \cite{amankwah2022qpixl} shows that the flexible representation of quantum images and related schemes reduce to a uniformly controlled rotation whose circuit, after a Walsh--Hadamard transform of the angles and Gray-code ordering of the controls, uses only $R_Y$ and CNOT gates and admits compression by discarding small transformed angles. \qcrank{} \cite{balewski2024qcrank} generalizes this to $n_d$ data qubits controlled by $n_a$ address qubits: a sequence of $N=2^{n_a}n_d$ real values $\theta_{a,d}$ is encoded by preparing the address register in uniform superposition and applying to each data qubit $d$ the uniformly controlled rotation
\begin{equation}\label{eq:bg:ucr}
U_d=\sum_{a=0}^{2^{n_a}-1}\ket{a}\!\bra{a}\otimes R_Y(\theta_{a,d}),\qquad
\ket{\Psi}=\frac{1}{\sqrt{2^{n_a}}}\sum_{a}\ket{a}\bigotimes_{d=1}^{n_d}\Bigl(\cos\tfrac{\theta_{a,d}}{2}\ket{0}+\sin\tfrac{\theta_{a,d}}{2}\ket{1}\Bigr).
\end{equation}
The M{\"o}tt{\"o}nen decomposition \cite{mottonen2004ucr} implements each $U_d$ with exactly $2^{n_a}$ single-qubit $R_Y$ rotations, whose transformed angles $\tilde{\bm{\theta}}=M^{-1}\bm{\theta}$ are obtained from a Walsh--Hadamard-type matrix $M$, interleaved with $2^{n_a}$ CNOTs whose controls follow a Gray code, as sketched in \cref{fig:bg:qcrank}. All $n_d$ data qubits are processed in parallel, so the total CNOT count is $\Theta(N)$ and the circuit depth is $\Theta(2^{n_a})$, independent of $n_d$. Readout is equally parallel: measuring all qubits and conditioning on the address bits yields $\expval{Z_d}_a=\cos\theta_{a,d}$, so every stored value is recovered from the same pool of shots as an expectation value rather than by tomography. With $S$ total shots each address is sampled about $S/2^{n_a}$ times, and the standard error of each recovered value scales as $\sqrt{2^{n_a}/S}$; this shot cost of reading out is the price of the parallelism and is quantified in the cost model of \cref{ch:synthesis} \cite{guo2026nonlinearwaves}. Balewski et al. demonstrated \qcrank{} on trapped-ion and transmon hardware \cite{balewski2024qcrank} and compiled it for a dynamically programmable neutral-atom processor \cite{balewski2025qcrankdpqa}. \qcrank{} circuits are a benchmark of \cref{ch:qgear}, the substrate of the attention and encoding primitives of \cref{ch:qml}, and the objects that \shardq{} cuts and reassembles.

% alt: QCrank circuit with two address qubits and one data qubit. Hadamard gates act on the address qubits; the data qubit receives an alternating sequence of four RY rotations with transformed angles and four CNOTs controlled by the address qubits in Gray-code order; all three qubits are measured at the end.
\begin{figure}[htbp]
\centering
\begin{quantikz}[row sep={0.75cm,between origins}, column sep=0.28cm]
\lstick{$a_0=\ket{0}$} & \gate{H} & & & \ctrl{2} & & & & \ctrl{2} & \meter{} \\
\lstick{$a_1=\ket{0}$} & \gate{H} & \ctrl{1} & & & & \ctrl{1} & & & \meter{} \\
\lstick{$d=\ket{0}$} & \gate{R_Y(\tilde{\theta}_0)} & \targ{} & \gate{R_Y(\tilde{\theta}_1)} & \targ{} & \gate{R_Y(\tilde{\theta}_2)} & \targ{} & \gate{R_Y(\tilde{\theta}_3)} & \targ{} & \meter{}
\end{quantikz}
\caption{\qcrank{} encoding of $N=4$ real values on $n_a=2$ address qubits and $n_d=1$ data qubit. The uniformly controlled rotation of \cref{eq:bg:ucr} is decomposed following \cite{mottonen2004ucr} into $2^{n_a}=4$ rotations by the transformed angles $\tilde{\theta}_j$ and four CNOTs whose controls follow a Gray code; additional data qubits repeat the same pattern in parallel. Schematic after \cite{balewski2024qcrank}.}
\label{fig:bg:qcrank}
\end{figure}

\subsection{Polynomial Computation, Block Encodings and Readout Cost}

Once real values are encoded as angles, arithmetic can be performed on them coherently. The EHands protocol of Balewski et al. \cite{balewski2025ehands} provides primitives for multiplication, weighted addition and polynomial evaluation on real-valued encoded states, so that a polynomial of the encoded value is produced without intermediate measurement; this is the starting point for the neural-native arithmetic of \cref{ch:qml} \cite{guo2026nnqa}, which compiles a trained polynomial network into such primitives. A different route to matrix arithmetic is block encoding: a unitary $U_A$ on $a+n$ qubits block-encodes an $n$-qubit operator $A$ with normalization $\alpha$ if $(\bra{0}^{\otimes a}\otimes I)\,U_A\,(\ket{0}^{\otimes a}\otimes I)=A/\alpha$. The FABLE compiler \cite{camps2022fable} produces approximate block encodings of dense matrices using the same uniformly controlled rotation structure as \qcrank{}, and quantum singular value transformation (QSVT) \cite{gilyen2019qsvt} then applies a polynomial transformation to the singular values of $A$ with a number of block-encoding calls equal to the polynomial degree, unifying Hamiltonian simulation, linear-system solving and amplitude amplification. The vectorized dot product of \cref{ch:qml} \cite{guo2025vqt} can be read as an approximate block encoding whose output is recovered as expectation values. In all these schemes the honest cost includes the shots needed to read the result: an expectation value estimated to precision $\eps$ requires $\bigO{1/\eps^2}$ shots, and when all $N$ outputs must be read, as in the nonlinear time stepping of \cref{ch:synthesis}, the readout cost scales as $\bigO{N/\eps^2}$ and can cancel the coherent speed-up \cite{guo2026nonlinearwaves}. Classical shadows \cite{huang2020shadows} reduce the cost of estimating many observables from randomized measurements and are the natural comparison point for such parallel readout.

\section{Quantum Machine Learning}\label{sec:bg:qml}

\subsection{Variational Classifiers and Transfer Learning}

Quantum machine learning (QML) \cite{biamonte2017qml} in the NISQ setting mostly takes the form of variational classifiers: a data-encoding circuit followed by a trainable circuit $U(\bm{\theta})$ and a measurement whose expectation value is the model output. Schuld et al. \cite{schuld2020circuitcentric} formalized this circuit-centric classifier and showed that its number of trainable parameters can be polylogarithmic in the input dimension when amplitude encoding is assumed, and Havl{\'\i}{\v{c}}ek et al. \cite{havlicek2019featurespaces} interpreted the encoding circuit as a feature map and ran quantum kernel and variational classifiers on superconducting hardware. Hybrid architectures interpose classical layers before and after the quantum circuit. Mari et al. \cite{mari2020transferlearning} introduced quantum transfer learning, in which a pretrained classical network is truncated and its features are compressed into a few qubits for a small variational classifier trained on a new task; this sequential dressed quantum circuit became the standard hybrid template. The \qpie{} network of \cref{ch:qml} \cite{guo2025qpie} departs from it by exchanging information between classical and quantum branches in parallel, and is evaluated on the same benchmark families (Hymenoptera, MNIST) so that the two can be compared directly.

\subsection{Gradients, Expressivity and Trainability}

Training requires gradients of $C(\bm{\theta})=\expval{O}{\psi(\bm{\theta})}$. On hardware, the parameter-shift rule \cite{crooks2019parametershift} evaluates the exact derivative with respect to the parameter of a gate $e^{-i\theta P/2}$, with $P$ a Pauli operator, from two shifted circuits,
\begin{equation}\label{eq:bg:paramshift}
\frac{\partial C}{\partial\theta}=\frac12\Bigl[C\bigl(\theta+\tfrac{\pi}{2}\bigr)-C\bigl(\theta-\tfrac{\pi}{2}\bigr)\Bigr],
\end{equation}
so that a full gradient costs $2|\bm{\theta}|$ circuit evaluations, each with its own shots. In simulation, adjoint differentiation \cite{jones2020adjoint} computes all gradient components in a single forward and backward sweep over the circuit at the memory cost of a few state-vector copies, which is why GPU simulators such as PennyLane Lightning are organized around it \cite{asadi2024pennylanelightning}. The dynamic gradient selection of \qpie{} \cite{guo2025qpie}, which uses parameter shift on QPUs and adjoint differentiation on GPUs, is a direct application of this distinction. Beyond gradient cost, the question of what a quantum model can represent is addressed through the Fisher information: Abbas et al. \cite{abbas2021power} used the spectrum of the Fisher information matrix to define an effective dimension and argued that well-designed quantum models have a flatter, less degenerate spectrum than comparable classical ones, which is the quantity \qpie{} reports. Expressivity is, however, in tension with trainability: the barren-plateau results of \cref{sec:bg:vqa} \cite{mcclean2018barren,holmes2022expressibility} apply verbatim to quantum neural networks, and the encoding-dependent analysis of Schuld et al. \cite{schuld2021dataencoding} shows that the data encoding, not only the trainable part, bounds the function class. This is one motivation for the gradient-free, structure-derived encodings of \cref{ch:qml}.

\subsection{Quantum Sequence Models}

The transformer \cite{vaswani2017attention} replaced recurrence with self-attention, whose core operation is the scaled dot product $\mathrm{softmax}(QK^{\mathsf{T}}/\sqrt{d})V$ over a sequence of length $T$, costing $\bigO{T^2d}$ per head. Quantum analogues have been proposed along two lines. Quantum recurrent models such as the quantum long short-term memory (Q-LSTM) of Chen et al. \cite{chen2022qlstm} replace the gates of an LSTM cell with variational circuits, while Quixer \cite{khatri2024quixer} builds a transformer-like model from linear combinations of unitaries and quantum singular value transformation and reports language-modelling perplexities on standard corpora; both remain variational and inherit the gradient and plateau costs discussed above. The vectorized quantum transformer of \cref{ch:qml} \cite{guo2025vqt} instead computes the attention matrix with a shot-efficient, gradient-free \qcrank{}-based dot product and is compared with Q-LSTM, Quixer and a classical nanoGPT baseline on the Brown corpus.

\section{Circuit Cutting and Knitting}\label{sec:bg:cutting}

Circuit cutting partitions a large circuit into sub-circuits small enough for the available hardware or simulator, executes them separately, and reconstructs the original circuit's output by classical post-processing, a step often called knitting. Two kinds of cut exist. A \emph{wire cut} severs a qubit wire between two gates; Peng et al. \cite{peng2020circuitcutting} showed that the identity channel on the cut wire can be expanded in a basis of measure-and-prepare operations, measuring in the Pauli bases and preparing Pauli eigenstates, so that the original expectation value becomes a signed sum of expectation values of sub-circuits. A \emph{gate cut} removes a two-qubit gate; Mitarai and Fujii \cite{mitarai2021virtualgate} showed that a CNOT or CZ can be written as a quasi-probability decomposition (QPD)
\begin{equation}\label{eq:bg:qpd}
\mathcal{U}=\sum_{i}c_i\,\mathcal{F}_i,\qquad c_i\in\R,\qquad \gamma=\sum_{i}|c_i|,
\end{equation}
into local single-qubit operations $\mathcal{F}_i$ with real coefficients that may be negative. Sampling the terms with probability $|c_i|/\gamma$ and weighting the outcomes by $\gamma\,\mathrm{sign}(c_i)$ yields an unbiased estimator whose variance is inflated by $\gamma^2$ relative to the uncut circuit. For a CNOT $\gamma=3$, so cutting $k$ gates independently incurs a sampling overhead of $\gamma^{2k}=3^{2k}$, and a wire cut incurs $4^2=16$ per cut. This exponential overhead is the fundamental limitation of the technique and the reason cut placement is an optimization problem in its own right. Bravyi, Smith and Smolin \cite{bravyi2016trading} analyzed the general trade between classical and quantum resources in such decompositions. Piveteau and Sutter \cite{piveteau2024knitting} showed that classical communication between the sub-circuits reduces the overhead of cutting $n$ parallel CNOTs from $\gamma=3^{n}$ to $\gamma=2^{n+1}-1$, and that gate cuts with communication are strictly more efficient than wire cuts. CutQC \cite{tang2021cutqc} was the first end-to-end system: it finds cuts with a mixed-integer program, runs the sub-circuits on small devices and reconstructs full probability distributions with dynamic-definition post-processing, demonstrating the evaluation of circuits larger than the device; IBM's qiskit-addon-cutting packages QPD-based gate and wire cutting for production use. The \shardq{} framework of \cref{ch:qml} \cite{guo2025shardq} specializes cutting to \qcrank{}-type data-encoding circuits: its SparseCut heuristic chooses gate cuts using physical qubit distance and calibrated error rates, and its reconstruction is global rather than pairwise, with the $\bigO{3^{2k}}$ overhead absorbed on GPUs.

\section{Quantum Error Correction and Fault Tolerance}\label{sec:bg:qec}

\subsection{Stabilizer Codes}

Quantum error correction (QEC) protects logical information by encoding it redundantly across many physical qubits so that errors can be detected and undone without measuring the logical state. The stabilizer formalism of Gottesman \cite{gottesman1997stabilizer} describes nearly all codes in practical use. Let $\mathcal{P}_n$ be the $n$-qubit Pauli group and $\mathcal{S}\subset\mathcal{P}_n$ an abelian subgroup not containing $-I$, generated by $n-k$ independent elements $S_1,\dots,S_{n-k}$. The code space is the joint $+1$ eigenspace
\begin{equation}\label{eq:bg:stabilizer}
\mathcal{C}=\bigl\{\ket{\psi}\in(\C^2)^{\otimes n}:\ S\ket{\psi}=\ket{\psi}\ \ \forall S\in\mathcal{S}\bigr\},\qquad \dim\mathcal{C}=2^{k},
\end{equation}
which encodes $k$ logical qubits. The logical operators are the elements of the normalizer $N(\mathcal{S})$ that are not in $\mathcal{S}$, and the code distance $d$ is the minimum weight of such an element; the code is written $\code{n}{k}{d}$ and corrects arbitrary errors on any $t=\lfloor(d-1)/2\rfloor$ qubits. A Pauli error $E$ that anticommutes with some generator flips the measured eigenvalue of that generator, and the vector of eigenvalues, the syndrome, identifies $E$ up to stabilizers; the general condition for a set of errors $\{E_a\}$ to be correctable is the Knill--Laflamme condition $PE_a^{\dagger}E_bP=c_{ab}P$ on the code projector $P$ \cite{knill1997theory}. The smallest code correcting an arbitrary single-qubit error is the $\code{5}{1}{3}$ perfect code of Laflamme et al. \cite{laflamme1996perfect}, which saturates the quantum Hamming bound: its $2^4=16$ syndromes exactly label the $3\times5=15$ weight-one Pauli errors plus the trivial one, a property exploited in \cref{ch:qec}, where all $3n$ weight-one faults are enumerated exhaustively. Steane's $\code{7}{1}{3}$ code \cite{steane1996codes} is a Calderbank--Shor--Steane (CSS) code \cite{calderbank1996goodcodes} built from the classical $[7,4,3]$ Hamming code, with separate $X$-type and $Z$-type generators; the CSS structure gives it transversal Clifford gates. Terhal \cite{terhal2015qecmemories} reviews the theory of QEC for quantum memories.

\subsection{Surface Codes, Bicycle Codes and Thresholds}

Kitaev's toric code \cite{kitaev2003anyons} and its planar variant, the surface code \cite{dennis2002topological,fowler2012surfacecodes}, place qubits on a two-dimensional lattice with weight-four $X$-type and $Z$-type checks on alternating faces, so that every check involves only nearest neighbours and syndrome extraction needs only local CNOTs and one ancilla per check. A distance-$d$ patch uses roughly $2d^2$ qubits. The threshold theorem states that below a physical error rate $p_{\mathrm{th}}$, increasing $d$ suppresses the logical error rate as $p_L\approx A\,(p/p_{\mathrm{th}})^{\lceil d/2\rceil}$; for the surface code under circuit-level depolarizing noise $p_{\mathrm{th}}$ is of order $1\%$ \cite{fowler2012surfacecodes}. Decoding maps a syndrome history to a likely error, usually by minimum-weight perfect matching, and Sundaresan et al. \cite{sundaresan2023subsystem} compared matching and maximum-likelihood decoders on a heavy-hex subsystem code over multiple rounds. The 2025 experiment of Google Quantum AI \cite{acharya2025belowthreshold} operated distance-3, -5 and -7 surface codes below threshold, with each increase of the distance by two reducing the logical error rate by a factor $\Lambda\approx2$. Because the rate $k/n$ of the surface code vanishes as $d$ grows, quantum low-density parity-check codes with higher rate are attractive: the bivariate bicycle codes of Bravyi et al. \cite{bravyi2024bicycle}, exemplified by the $\code{144}{12}{12}$ code with weight-six checks, promise roughly a tenfold reduction in overhead relative to the surface code at comparable noise, at the cost of non-planar connectivity. Syndrome extraction circuits are themselves faulty, and a fault on an ancilla can propagate to several data qubits, which fault-tolerant extraction schedules and repeated measurement rounds are designed to contain. Instead of correcting, many small experiments post-select: runs with a nontrivial syndrome are rejected, which lowers the error rate of the accepted runs at the cost of discarded shots. The security analysis of \cref{ch:qec} \cite{guo2026qecaudit} is built precisely around what survives such rejection.

\subsection{Clifford Simulation, Random Ensembles and Magic}

The Clifford group $\mathcal{C}_n$ is the normalizer of the Pauli group and is generated by $H$, $S$ and CNOT. The Gottesman--Knill theorem states that circuits of Clifford gates, Pauli measurements and stabilizer-state inputs can be simulated efficiently on a classical computer, because the state can be tracked by its $n$ stabilizer generators rather than by $2^n$ amplitudes; Aaronson and Gottesman \cite{aaronson2004stabilizer} gave a tableau representation with $\bigO{n}$ cost per gate and $\bigO{n^2}$ per measurement. Stim \cite{gidney2021stim} implements this with bit-packed, vectorized tableaus and a fast Pauli-frame sampler for repeated noisy runs, which allows syndrome-extraction circuits with thousands of qubits to be sampled and makes it the standard tool for estimating logical error rates and thresholds; the gate-level simulations of \cref{ch:qec} at $n=17$ use it. Encoding circuits of stabilizer codes are Clifford, which has two consequences exploited in this dissertation. First, ensembles of random Clifford unitaries are cheap to sample and to simulate, whereas Haar-random unitaries on $n$ qubits require exponentially many gates. Second, the Clifford group is an exact unitary 2-design \cite{dankert2009designs}, and in fact a 3-design \cite{webb2016clifford3design}: for every operator $X$ on two copies of the system,
\begin{equation}\label{eq:bg:design}
\mathop{\mathbb{E}}_{U\sim\mathcal{C}_n}\bigl[U^{\otimes2}X(U^{\dagger})^{\otimes2}\bigr]=\mathop{\mathbb{E}}_{U\sim\mathrm{Haar}}\bigl[U^{\otimes2}X(U^{\dagger})^{\otimes2}\bigr],
\end{equation}
so that any quantity depending on at most the second moment of the ensemble, including average fidelities and the disturbance measures of \cref{ch:qec}, takes its Haar value when averaged over Cliffords; classical shadows \cite{huang2020shadows} rely on the same fact. Clifford gates alone are not universal, and the missing ingredient, often called magic, is supplied by the $T$ gate. In the surface code $T$ cannot be implemented transversally and is instead injected through magic-state distillation \cite{bravyi2005magic}, whose cost dominates fault-tolerant resource estimates; the $T$-count is therefore the primary cost driver of a fault-tolerant circuit, ahead of Clifford count and depth, which is why it is a weighted term in the synthesis rubric of \cref{ch:synthesis} \cite{guo2026rubriq} and why the cryptanalytic estimates of \cref{sec:bg:crypto} are dominated by $T$-consuming arithmetic.

\section{Quantum Threats, Post-Quantum Cryptography and Standards}\label{sec:bg:crypto}

\subsection{Cryptanalytic Quantum Algorithms and Resource Estimates}

Shor's algorithm \cite{shor1997factoring} factors integers and computes discrete logarithms in polynomial time by using the quantum Fourier transform to find the period of modular exponentiation, which breaks RSA, finite-field Diffie--Hellman and elliptic-curve cryptography. Grover's algorithm \cite{grover1996search} finds a marked item among $N$ with $\bigO{\sqrt{N}}$ queries and thereby halves the effective key length of symmetric ciphers and hash functions; the standard response is to double key sizes rather than to change algorithms. The gap between these asymptotic statements and an actual attack is set by resource estimates. Gidney and Eker{\aa} \cite{gidney2021rsa} estimated that factoring a 2048-bit RSA modulus would take about eight hours on 20 million noisy superconducting qubits protected by a surface code at a physical error rate of $10^{-3}$, with the $T$-gate-consuming modular arithmetic dominating; Gidney's 2025 update \cite{gidney2025million} lowers the estimate to under a million noisy qubits and roughly a week of runtime by improving the arithmetic and the error-correction layout\todo{verify the qubit count and runtime of the 2025 estimate against the preprint}. These estimates connect the fault-tolerant cost accounting of \cref{sec:bg:qec} to concrete policy timelines. Mosca's inequality \cite{mosca2018cybersecurity} states that if the time $x$ for which data must remain secret plus the time $y$ needed to migrate infrastructure exceeds the time $z$ until a cryptographically relevant quantum computer exists, then the data is already at risk, because an adversary can harvest encrypted traffic now and decrypt it later.

\subsection{Standardization and Validation}

Post-quantum cryptography (PQC) \cite{bernstein2017pqc} replaces number-theoretic assumptions with problems believed to be hard for quantum computers, chiefly structured lattices, hash functions and error-correcting codes. NIST opened its standardization process with a public call for algorithms in December 2016 \cite{nist2016pqccall}; after three rounds of public cryptanalysis it announced the selected algorithms in July 2022 \cite{nist2022ir8413}, and in August 2024 it published three Federal Information Processing Standards: FIPS 203, the module-lattice key-encapsulation mechanism ML-KEM derived from CRYSTALS-Kyber \cite{nist2024fips203}; FIPS 204, the module-lattice digital signature ML-DSA derived from CRYSTALS-Dilithium \cite{nist2024fips204}; and FIPS 205, the stateless hash-based signature SLH-DSA derived from SPHINCS$^{+}$ \cite{nist2024fips205}. A fourth signature standard based on FALCON and an additional code-based key-encapsulation mechanism, HQC, are in progress. Standardizing an algorithm is not the same as trusting an implementation. FIPS 140-3 \cite{nist2019fips1403}, aligned with ISO/IEC 19790, defines four security levels for cryptographic modules and the Cryptographic Module Validation Program under which an accredited laboratory tests a module's approved algorithms, key management, self-tests and physical security before it may be used in federal systems. Among the approved components, deterministic random bit generators (DRBGs) are specified in SP 800-90A \cite{nist2015sp80090a}, which defines Hash\_DRBG, HMAC\_DRBG and CTR\_DRBG together with their instantiation, reseeding and health-test requirements; randomness quality is a validation item in its own right because every PQC key and signature depends on it. The Light Rider QPC SDK of \cref{ch:pqc} \cite{guo2026lightriderqpc} is a FIPS 140-3-style PQC module, and the randomness audit of \cref{ch:qec} \cite{guo2026qecaudit} deliberately borrows the seeding and reseeding vocabulary of SP 800-90A.

\subsection{Threat Models for Cloud Quantum Hardware}

Quantum key distribution (QKD) addresses the key-exchange problem physically rather than computationally. BB84 \cite{bennett1984bb84} encodes key bits in non-orthogonal photon states so that eavesdropping disturbs them detectably; device-independent QKD \cite{vazirani2014diqkd} certifies security from Bell-inequality violations without trusting the devices; and wiretap-channel analyses such as \cite{pan2020secretkey} bound the secret-key rate under restricted eavesdropping. QKD secures links rather than computations and is complementary to PQC. A newer question is the security of computation performed on shared, cloud-accessed QPUs, where a user's circuit runs on hardware the user neither owns nor observes. Blind quantum computation \cite{broadbent2009blind} hides the computation from the server, and Mahadev's protocol \cite{mahadev2018verification} lets a classical client verify a quantum server's output under cryptographic assumptions, but neither is practical on current devices. At the hardware level, multi-programming, or co-tenancy, executes circuits of different users simultaneously on one chip, and Ash-Saki et al. \cite{ashsaki2020crosstalk} showed that a co-tenant can exploit crosstalk to degrade or bias a victim's circuit, establishing fault injection through shared physical resources as a threat vector. The audit of \cref{ch:qec} \cite{guo2026qecaudit} formalizes such an adversary under a bounded fault model and asks how the public, fixed structure of QEC encoders exposes them, an emerging topic at the intersection of error correction and systems security.

\section{Large Language Models for Code and Circuit Synthesis}\label{sec:bg:llm}

Large language models (LLMs) trained on source code generate syntactically valid programs from natural-language or partial-code prompts. Codex \cite{chen2021codex} established the functional-correctness evaluation methodology of the field, in which generated programs are executed against held-out tests rather than compared textually to a reference, and open models such as Qwen2.5-Coder \cite{hui2024qwen25coder} now reach strong performance at 7--32 billion parameters. Quantum circuits expressed in \qiskit{}, OpenQASM or \cudaq{} are programs, so the same models can be prompted to synthesize circuits. Circuit synthesis is in one respect an unusually favourable target for such models, because a classical simulator can score a generated circuit exactly against its specification: the fidelity to a target state or unitary, adherence to a coupling graph and gate counts such as the $T$-count of \cref{sec:bg:qec} are all computable, and the simulator therefore doubles as an oracle. Supervised fine-tuning alone, however, does not teach a model to respect such quantitative constraints, and generated code that fails to parse or to compile is common.

Reinforcement learning from verifiable rewards supplies the missing feedback. Proximal policy optimization \cite{schulman2017ppo} maximizes a clipped surrogate objective and requires a learned value network; group relative policy optimization (GRPO), introduced with DeepSeekMath \cite{shao2024deepseekmath}, removes the value network by sampling a group of $N$ completions per prompt and using the group-normalized reward $\hat{A}_i=(r_i-\bar{r})/\sigma_r$ as the advantage, and DeepSeek-R1 \cite{deepseek2025r1} showed that GRPO with rule-based rewards can induce long-form reasoning. Fine-tuning a multi-billion-parameter model is made affordable by low-rank adaptation (LoRA) \cite{hu2022lora}, which freezes the pretrained weights $W_0$ and learns a low-rank update $W=W_0+BA$ with $\mathrm{rank}(BA)=r\ll\min(d_{\mathrm{in}},d_{\mathrm{out}})$, reducing the number of trainable parameters by orders of magnitude, and by DeepSpeed ZeRO \cite{rajbhandari2020zero}, which partitions optimizer states (stage~1), gradients (stage~2) and parameters (stage~3) across data-parallel GPUs so that the memory footprint of a training run shrinks with the number of devices. Generative alternatives to LLMs exist: F{\"u}rrutter et al. \cite{furrutter2024genqc} synthesize circuits with diffusion models conditioned on the target unitary. The \rubriq{} system of \cref{ch:synthesis} \cite{guo2026rubriq} combines these ingredients, a 7-billion-parameter coder with rank-64 LoRA, GRPO with a group size of eight and ZeRO stage 2 across A100 nodes, with a programmatic rubric reward evaluated by \cudaq{} simulation inside the loop, which places the GPU simulator of \cref{sec:bg:sim} at the centre of the learning process.

\section{Positioning of This Dissertation}\label{sec:bg:summary}

The preceding sections identify, for each layer of the stack in \cref{fig:bg:stack}, a body of established results and an unresolved gap. We close by stating those gaps in terms of the five research questions of \cref{ch:intro}.

\textbf{RQ1 (Scale).} GPU state-vector simulators and distributed frameworks exist (\cref{sec:bg:sim}), and \cudaq{} in particular reaches thousand-GPU scale, but each requires programs to be written in its own model. \qiskit{}-native workflows, which dominate algorithm development and hardware access, either run on Aer with limited multi-node scaling or must be ported by hand, and no published pipeline takes \qiskit{} circuits unchanged to distributed GPU simulation with containerized, Slurm-ready deployment or characterizes where communication overtakes computation for typical algorithm circuits at 30--42 qubits on a production HPC system. \Cref{ch:qgear} addresses this gap with \qgear{} \cite{guo2025qgear,guo2025qgearsoftware}.

\textbf{RQ2 (Noise).} QAOA on hardware is limited by the depth budget of \cref{sec:bg:nisq} and by the optimization cost and barren plateaus of \cref{sec:bg:vqa}. Warm starts and parameter concentration reduce but do not eliminate the black-box loop, and ZNE is usually applied as a generic post-processing step disconnected from the ansatz. Missing is a formulation in which the problem coefficients determine the angles directly, the ansatz is built from the device's native entangling gate, and noise amplification is designed jointly with the circuit, together with an oracle whose outputs are interpretable without tomography. \Cref{ch:nisq} addresses this gap with \deal{} \cite{guo2025deal} and \qawa{} \cite{guo2025qawa}.

\textbf{RQ3 (Data).} The encodings of \cref{sec:bg:encoding} make loading and parallel readout efficient, and the models of \cref{sec:bg:qml} are learnable, but the two bodies of work have rarely been combined. \qcrank{}-style vectorized encodings had not been made learnable, used as attention primitives, cut hardware-aware into sub-circuits, or used as the target of an exact compilation of a trained network; gradient-based hybrid models, conversely, were evaluated without accounting for shots or for the QPU-versus-GPU gradient trade-off. \Cref{ch:qml} addresses this gap with \qpie{} \cite{guo2025qpie}, \vqt{} \cite{guo2025vqt}, \shardq{} \cite{guo2025shardq} and \nnqa{} \cite{guo2026nnqa}.

\textbf{RQ4 (Synthesis).} LLM- and RL-based program synthesis (\cref{sec:bg:llm}) had not been coupled to the fault-tolerant cost metrics of \cref{sec:bg:qec} through a simulator-in-the-loop reward, and diffusion-based circuit synthesis targets unitaries rather than constraint-aware programs. At the same time, quantum algorithms for nonlinear partial differential equations hide the measurement-and-reload cost of \cref{sec:bg:encoding} inside linear embeddings \cite{jin2024schrodingerization}, so that no end-to-end account of when a quantum solver pays off exists. \Cref{ch:synthesis} addresses this gap with \rubriq{} \cite{guo2026rubriq,guo2026rubriqsoftware} and with the split-step cost model for nonlinear waves \cite{guo2026nonlinearwaves,guo2026splitstepqude}.

\textbf{RQ5 (Trust).} QEC theory (\cref{sec:bg:qec}) assumes an adversary-free noise model, and the cloud threat models of \cref{sec:bg:crypto} have been studied for NISQ circuits but not for encoders and syndrome-based rejection; nor had QEC been connected to the randomness and module-validation standards that govern classical cryptography. \Cref{ch:qec} addresses this gap with the structured-randomness audit \cite{guo2026qecaudit}, and \cref{ch:pqc} with a FIPS 140-3-style post-quantum module \cite{guo2026lightriderqpc} positioned against the ML-KEM, ML-DSA and SLH-DSA standards \cite{nist2024fips203,nist2024fips204,nist2024fips205}, FIPS 140-3 \cite{nist2019fips1403} and SP 800-90A \cite{nist2015sp80090a}.

Together, these five gaps trace the outline of \cref{fig:bg:stack}: a QPU is useful only to the extent that the classical layers around it are fast, structured, data-efficient, cost-aware and trustworthy, and the chapters that follow build those layers in turn.
